% Apertura conceptual de la Parte I. Esta sección reúne, en una sola
% residencia, el problema histórico del continuo, la distinción entre marca,
% intervalo y medida, y la descomposición nonádica que conduce a APP.

\section{Le problème du continu et la généalogie de la mesure}
\label{sec:continuo-numero-medida-hmt-revfinal}
\index{continu!structure discrète}
\index{nombre!généalogie}
\index{mesure!couplage}

Cette ouverture développe la dimension nombre--mesure de la
\hyperref[sec:tesis-inaugural-helicoidal-hmt-md]{thèse inaugurale du traité}:
la structure discrète construit le continu et la carte réelle publie ensuite
sa coordonnée, sans intervenir comme domaine générateur.

L'histoire du continu peut se lire comme l'histoire d'une séparation et
de ses réconciliations successives. L'arithmétique antique organisait les
pluralités ; la géométrie comparait les longueurs, les aires et les volumes. La théorie
grecque des proportions a permis de raisonner sur des grandeurs incommensurables :
la diagonale du carré unité existait comme construction exacte alors même que
sa longueur ne pouvait pas se publier comme rapport d'entiers. Les livres V et VII
des \emph{Éléments} ont ainsi conservé deux domaines, grandeur et nombre,
unis par des règles de comparaison mais dotés d'opérations distinctes
\cite{EuclidElements,SEPAristotleMath}.

L'analyse moderne a donné à cette réconciliation une forme décisive. Les
suites de Cauchy, les coupures de Dedekind et la théorie des ensembles ont fait
de \(\mathbb R\) un corps ordonné complet, capable de publier dans un même
langage arithmétique les grandeurs continues. Cantor a distingué cardinalité
et ordre ; Dedekind a fait de la continuité une propriété structurelle du
système numérique ; Weierstrass a établi la discipline locale de la limite et de
l'approximation \cite{SEPDedekind,SEPContinuity,SEPSetTheory}. Le continu réel
qui en résulte est l'une des constructions les plus puissantes de l'histoire des
mathématiques : il identifie les présentations qui convergent vers la même valeur et rend
possibles le calcul, la géométrie analytique et la formulation quantitative des
théories physiques.

L'holographie modulaire triadique pose une question antérieure à cette
identification : \emph{quelle structure finie et prolongeable produit la valeur
que publie la droite réelle ?} Son objet premier est la généalogie compatible d'une
évaluation. Un chiffre est la sortie d'un lecteur positionnel ; une valeur est
la coordonnée scalaire obtenue en contractant une famille de cylindres ; un nombre
HMT est la section qui conserve, en outre, la phase, l'orientation, le feuillet
arithmétique, le quotient d'élévation, la retenue, la frontière et la mémoire de
la route. L'évaluation archimédienne prend ainsi la forme
\[
 \operatorname{ev}_{\mathbb R}:X_\infty^{\rm HMT}\longrightarrow\mathbb R,
\]
et ses fibres réunissent des histoires qui partagent une coordonnée sans perdre pour autant
leur identité opératoire.

Cette formulation change l'ordre explicatif de la mesure. Compter détermine
la cardinalité d'un ensemble fini ; mesurer couple un état, un lecteur et
une carte, restreint leurs continuations compatibles et publie une observable.
Si \(X\) est l'espace d'états, \(A\) l'appareil ou le lecteur et
\(\mathcal C\subset X\times A\) la condition de compatibilité, le schéma
minimal est
\[
 (x,a)\in\mathcal C
 \longmapsto \mathcal O(x,a)
 \longmapsto \mathcal U\bigl(\mathcal O(x,a)\bigr),
\]
où \(\mathcal O\) conserve le type de l'observable et \(\mathcal U\)
l'exprime dans la carte choisie. Le chiffre final constitue la coordonnée visible
du couplage complet. Cette distinction sera la base mathématique de la
théorie de l'observateur et de la mécanique dimensionnelle dans la Partie V.

La conséquence épistémologique est précise. La métrologie physique ordinaire
attribue des coordonnées aux observables et les théories relient ensuite ces
coordonnées par des équations ; HMT construit d'abord l'état, la route et le
lecteur dont descend la coordonnée. Les constantes fondamentales occupent
alors la place de résultats d'une généalogie et la confrontation expérimentale
s'applique à la sortie. Dans le même mouvement, une symétrie HMT exprime la
conservation de l'unité à travers une transformation : la valeur enregistre la
coordonnée publiée et la projection exceptionnelle enregistre l'organisation qui
a survécu au transport. La structure et le nombre deviennent deux
aspects vérifiables d'une même évolution discrète.

\subsection{Marques, intervalles et cellule nonadique}

La construction finie commence sur la droite des entiers. On fixe
l'origine \(0\), le seuil décimal \(10\) et la cellule des représentants
positifs
\[
 E_9=\{1,2,3,4,5,6,7,8,9\}.
\]
Les éléments de \(E_9\) remplissent deux fonctions liées et distinctes :
ce sont les neuf marques de phase du chemin \(P_9\) et ce sont aussi les neuf
coordonnées entières positives \(d(0,r)=r\) mesurées depuis l'origine. Cette
seconde lecture n'ajoute pas d'arêtes au chemin. Les neuf marques déterminent
exactement les huit intervalles intérieurs semi-ouverts
\[
 J_r=[r,r+1),\qquad 1\le r\le8.
\]
L'arête de fermeture est représentée dans la carte relevée par \([9,10)\) ; sa
projection nonadique revient à la première phase et conserve la retenue du
tour suivant. L'origine \(0\) appartient uniquement à la carte décimale étendue
et n'ajoute pas d'intervalle intérieur à \(P_9\). Marques, coordonnées, intervalles et
fermeture sont ainsi typés avant d'appliquer une quelconque opération.

Pour chaque \(n\ge1\), la division euclidienne nonadique définit
\[
 \rho_9(n)=1+((n-1)\bmod9),
 \qquad
 q_9(n)=\left\lfloor\frac{n-1}{9}\right\rfloor,
\]
\begin{equation}
 n=9q_9(n)+\rho_9(n).
 \label{eq:rueda-helice-nonadica-revfinal}
\end{equation}
La coordonnée \(\rho_9\) donne la position sur une roue à neuf phases ;
\(q_9\) conserve la hauteur accumulée. Le successeur agit par
\[
 (q,r)\longmapsto
 \begin{cases}
 (q,r+1),&1\le r<9,\\
 (q+1,1),&r=9.
 \end{cases}
\]
Ainsi, \(n\) et \(n+9\) partagent une phase et occupent des hauteurs consécutives. La
réalisation géométrique
\[
 \mathcal H_9(n)=
 \left(
  \cos\frac{2\pi(\rho_9(n)-1)}9,
  \sin\frac{2\pi(\rho_9(n)-1)}9,
  q_9(n)+\frac{\rho_9(n)-1}{9}
 \right)
\]
rend cette structure visible : projetée sur le plan, c'est une roue ;
conservée dans les trois coordonnées, c'est une hélice. Le retour \(9\to1\)
récupère l'angle et élève l'histoire. C'est la forme élémentaire de la
distinction qui régira tout le TPK :
\[
 \boxed{\text{retour de phase}\;\neq\;\text{retour de l’état}.}
\]

\subsection{D'une cellule à deux coordonnées : naissance d'APP}

Le produit cartésien de deux cellules nonadiques construit une carte
bidimensionnelle. Son noyau intérieur contient \(8^2=64\) positions. La neuvième
ligne et la neuvième colonne publient la classe résiduelle nulle au moyen de l'étiquette
\(9\) et forment un gnomon de dix-sept cellules :
\[
 81=64+(9+8)=64+17.
\]
L'identification modulaire des côtés opposés transforme la carte carrée
en le tore
\[
 X_9=(\mathbb Z/9\mathbb Z)^2,
\]
et les translations cardinales le dotent d'un graphe de Cayley. L'opération
géométrique est, par conséquent, une séquence précise : noyau
\(8\times8\), complétion gnomonique \(9\times9\), identification toroïdale et
transport par le graphe.

Sur chaque sommet \((i,j)\in X_9\) agissent deux évaluations du même support :
le feuillet additif \(\Sigma\) et le feuillet multiplicatif \(\Pi\). Leur
incompatibilité exacte apporte une dynamique, car une route doit conserver quel
feuillet a agi, quel résidu il a publié et quel quotient il a produit. APP construit le
support et ses deux évaluations ; TRIT oriente localement la transition ; TPK
transporte l'état complet, conserve simultanément les deux feuillets et, au moyen de
\(M_{\rm ph}\), type le mode et la phase TRIT de la transition ; il conserve en outre la
mémoire de l'élévation et décide de la survie des prolongations.
L'architecture inaugurale est fixée par la chaîne
\[
 \boxed{
 E_9\longrightarrow
 \mathrm{APP}\longrightarrow
 \mathrm{TRIT}\longrightarrow
 \mathrm{TPK}.}
\]

L'importance de ce commencement réside dans sa capacité de prolongation. La
même différence entre position et hauteur qui transforme la roue en hélice se
reproduit dans la périodicité de \(108=12\cdot9\) itérations, dans le système
cyclique orienté à douze phases, dans
\hyperref[sec:umbrales-prolongacion-tpk-inaugural]{l'architecture de niveaux et de
seuils de la prolongation TPK}
\(w_6\to w_{12}\to w_{18}\to w_{24}\to w_{30}\to R_{36}\to G_9\), dans la dérivation de
\(\pi,e,\varphi\) et de \(\alpha_{\rm HMT}\), et dans la double projection vers la
valeur réelle et l'incidence exceptionnelle. La structure discrète du continu
désigne précisément cette continuité d'une généalogie finie à travers des
profondeurs arbitraires.

\begin{statusbox}
L'histoire du nombre, de la grandeur et du continu appartient au contexte mathématique
classique. La décomposition \eqref{eq:rueda-helice-nonadica-revfinal}, le tore
nonadique, les feuillets d'APP et le transport de mémoire sont des résultats du
corpus HMT. L'immersion hélicoïdale \(\mathcal H_9\) est une formalisation
géométrique de cette décomposition. L'évaluation sur tout \(\mathbb R\)
est attestée par le théorème de couverture et de fermeture archimédienne du
continu généalogique
\cref{thm:espiral-cierre-arquimediano-continuo-genealogico} ; la généalogie des
valeurs construites se formule dans le système inverse des sections
compatibles.
\end{statusbox}
